
% =====================================================================
\section{T6: histories do not return}
\label{sec:monotone}
% =====================================================================

A gated trajectory may revisit a vertex; we show it can never return to a prior
\emph{state}, because the committed count---the number of steps taken---is
strictly monotone and is part of the state. Revisiting is not returning;
resemblance of position is not identity of state. This is the single
order-theoretic law that the framework instantiates wherever ``the same thing
again'' is claimed.

\subsection{The committed count is strictly monotone}

\begin{definition}[State of a walk]
\label{def:state}
The \emph{state} of a walk $\Walk=(v_0,\dots,v_\ell)$ at step $\ell$ is the pair
\[
s_\ell\;=\;\bigl(v_\ell,\ \rec(\Walk)\bigr),\qquad \rec(\Walk)=\ell,
\]
the current vertex together with the committed count of steps taken to reach it.
Two states are equal iff both components agree.
\end{definition}

\begin{theorem}[Non-return]
\label{thm:nonreturn}
Along any walk in a contact graph the committed count is strictly increasing:
$\rec$ after step $i+1$ exceeds $\rec$ after step $i$ by exactly one. Consequently
no walk returns to a previously held state: for $i<j$, $s_i\neq s_j$ even when
$v_i=v_j$. A step intended to undo a previous step is itself a step and increases
$\rec$. Hence resemblance of position \textup{(}$v_i=v_j$\textup{)} never entails
identity of state \textup{(}$s_i=s_j$\textup{)}.
\end{theorem}

\begin{proof}
Each step appends one edge to the committed multiset and advances the index by
one, so $\rec$ increases by exactly $1$ per step; it is therefore strictly
increasing and, for $i<j$, $\rec_i=i<j=\rec_j$. Thus the second components of
$s_i$ and $s_j$ differ, so $s_i\neq s_j$ regardless of whether $v_i=v_j$. An
``undo'' move from $v_j$ back to $v_{j-1}$ is a further step, producing state
$(v_{j-1},j{+}1)\neq(v_{j-1},j{-}1)=s_{j-1}$; it does not restore the earlier
state but creates a new one with higher count.
\end{proof}

\begin{corollary}[Cycles in position are not cycles in state]
\label{cor:cycles}
If a walk traverses a closed walk in $\Gph$ \textup{(}$v_0=v_\ell$, $\ell>0$\textup{)},
its state has nonetheless advanced: $s_0=(v_0,0)\neq(v_0,\ell)=s_\ell$. The graph
may be revisited arbitrarily often; the state never repeats.
\end{corollary}

\begin{proof}
Immediate from \Cref{thm:nonreturn} with $i=0$, $j=\ell$, $v_0=v_\ell$.
\end{proof}

\subsection{Re-contact is forward only}

The next result is the form non-return takes for ``returning to'' an earlier
vertex deliberately---re-contact. It is reached forward, at a higher count, hence
is a new state, not the old one.

\begin{proposition}[Forward re-contact]
\label{prop:recontact}
Let a walk reach vertex $v$ at step $i$ and, later, reach $v$ again at step
$j>i$ \textup{(}a deliberate re-contact of $v$\textup{)}. Then the re-contact
state $s_j=(v,j)$ is strictly later than the original $s_i=(v,i)$ on the
committed order, and $s_j\neq s_i$. There is no operation that re-contacts $v$ at
its original state; every re-contact is at a state in the future of the original.
\end{proposition}

\begin{proof}
By \Cref{thm:nonreturn}, $\rec$ is strictly increasing, so $j>i$ forces
$s_j\neq s_i$; the re-contact carries the count $j>i$, placing it later on the
order. No step decreases $\rec$, so no re-contact can carry the original count
$i$ again.
\end{proof}

\subsection{The progressing whole: non-return from a growing complement}

The committed count gives non-return internally. There is a second, independent
source of non-return when the whole itself progresses: the complement against
which a part is individuated grows, so a part is individuated against a strictly
larger negation-set at each later stage. We record this, as it grounds the
forward-only character in the non-completability of \Cref{sec:ground}.

\begin{definition}[Progressing whole]
\label{def:progressing}
A whole $\med$ \emph{progresses} if its committed history grows: at stage $j>i$
the set of parts and committed boundaries present includes those present at stage
$i$ and at least one more. Equivalently, the medium relative to a fixed part
$A$---its complement $\med\setminus A$---is non-decreasing and strictly larger at
some later stage.
\end{definition}

\begin{theorem}[Temporal non-return]
\label{thm:temporal-nonreturn}
Let $\med$ be a non-completable, progressing whole and $A$ a fixed thing. Then
$A$ at stage $j>i$ is individuated against a strictly larger complement than $A$
at stage $i$; hence the individuation-state of $A$ at $j$ differs from that at
$i$, even though $A$'s interior residue \textup{(}its conserved identity\textup{)}
is unchanged. A thing is therefore distinguishable from itself at every earlier
and later stage: its conserved interior the same, its individuating exterior never
the same.
\end{theorem}

\begin{proof}
By \Cref{def:progressing} the complement $\med\setminus A$ at stage $j$ properly
contains that at stage $i$. Individuation of $A$ is determination against its
complement (\Cref{thm:negation}); a strictly larger complement is a different
negation-set, so the individuation-state---the pair of $A$ with the complement it
is individuated against---differs between $i$ and $j$. The interior residue of $A$
is its private invariant (\Cref{def:private}), conserved under reshuffling
(\Cref{thm:private}), hence unchanged. Thus $A$-at-$j$ and $A$-at-$i$ share the
conserved interior but differ in the individuating exterior, and are
distinguishable as individuation-states while identical in conserved identity.
\end{proof}

\begin{corollary}[Same interior, different state, across all stages]
\label{cor:same-different}
For a fixed thing in a progressing whole, ``the same thing'' across stages is a
resemblance of conserved interior against an ever-growing exterior. By
\Cref{thm:nonreturn} and \Cref{thm:temporal-nonreturn} this resemblance is never
identity of state: the committed record grows and the complement grows, so no
later stage returns to an earlier individuation-state.
\end{corollary}

\begin{proof}
Combine the internal monotone count (\Cref{thm:nonreturn}) with the growing
complement (\Cref{thm:temporal-nonreturn}); both strictly advance the
individuation-state, while the interior residue (\Cref{def:private}) is conserved.
\end{proof}

\begin{remark}[One law, several guises---offered as orientation]
\label{rem:onelaw}
\Cref{thm:nonreturn} is a single order-theoretic fact: in a finite system whose
acts are counted, the count is monotone, so a revisited configuration is a new
state. A reader may recognise this as the common form of several apparently
distinct impossibilities---that a reversed process is not the initial one, that a
re-contacted memory is not the original experience, that a reconstituted system
bearing the mark ``was interrupted'' is individuated by a negation the original
lacked. Each is ``resemblance under a monotone record is not identity,'' the
content of \Cref{thm:nonreturn} and \Cref{prop:recontact}. We state and use only
the graph fact; the guises are noted, not relied upon.
\end{remark}

\begin{remark}[Why finiteness matters here too]
\label{rem:nonreturn-finite}
The count $\rec$ is well-defined and strictly monotone because steps are
discrete and a walk's length is a natural number---features of a finite graph
traversed in discrete steps. The law is not about any continuous parameter; it is
the order on committed steps. This keeps T6 within the same premise as the rest:
finite parts, discrete separations, a counted history.
\end{remark}
